\section{Continuing the imbalance sum}\label{sec:E-interpolation}

The number of reference particles must next be varied with \(M\).
There is no integral of noninteger dimension to which Cauchy's theorem
could be applied directly. We construct a holomorphic approximation
from polynomial moments of the integer-dimensional integrals. The
construction has two outputs: arbitrarily small absolute error on the
integer grid, and a coarse growth bound on a complex dimension polydisk.
The grid-transfer lemma at the end of the section uses both outputs to
carry the sharp integer-dimension estimates to the complex dimensions
needed for the imbalance contour shift.

\begin{proposition}[Interpolation of the reference dimension]\label{prop:04-interpolation}
Fix polynomial bounds on \(n_e,\ell_e,|M|\). For every fixed \(b>0\)
there is a jointly holomorphic function \(\widehat F(\nu,M)\), with
\(\nu=(\nu_+,\nu_-)\in\mathbb C^2\), on a neighborhood of the indicated \(M\)-disk and the polydisk
\(|\nu_\sigma-m|\le\sqrt m\), such that
\begin{equation}
 |\widehat F(\nu,M)|\le \exp(L^C),\qquad
 |\widehat F(\nu,M)-F_\nu(M)|\le e^{-L^b}
       \quad\text{for }\nu\in\mathbb Z_{\ge0}^2\text{ in the polydisk}.
 \label{eq:G-20}
\end{equation}
Here \(C\) is fixed after the truncation bounds and \(b\) are fixed.
The bounds are uniform in the real explicit-particle positions on their
labeled lines. The integrals include their reference masses
\(J_{\nu_+}(w_0)J_{\nu_-}(w_0)\), but no external phase factor.
\end{proposition}

The next three subsections prepare the proof of
Proposition~\ref{prop:04-interpolation}: polynomial approximation,
a finite endpoint-subset expansion, and continuation of the retained
factorial moments. The proof is assembled after these constructions.

\subsection{Polynomial approximation under the integer laws}

We first make all errors that will be estimated under a positive
integer-dimensional probability law. In \(\Phi_M\), truncate the
physical sine expansions of the linear dual test and of the primitive
kernel \(\mathscr K\). The degrees can be polynomial in \(L\), with
error \(e^{-L^{B_0}}\) in any prescribed fixed smooth norm. Indeed a
Gevrey bound \(A^k(k!)^P\) and \(k\) integrations by parts bound the
\(n\)-th Fourier coefficient by \(C A^k(k!)^P n^{-k}\);
optimizing \(k\) gives \(C\exp[-c(n/A)^{1/P}]\). The Gevrey scale
here is polynomial in \(L\). Fixed Fourier cutoffs preserve
holomorphy in \(M\).

The total variation of each particle or mean measure is
\(O(e^{\gamma L})\), so the exponent error is bounded by
\(e^{-L^{B_0}+O(L)}\) uniformly on configurations. Independently of
this Fourier cutoff, Lemma~\ref{lem:04-quadratic} bounds the absolute
expectation of the original exponential by \(\exp(L^{C_0})\), using
the gap, polynomial norm of the linear term, and polynomial lower
bound on the real constant. Since
\(|e^{-z}-1|\le |z|e^{|z|}\), choosing \(B_0\) larger than the
required precision and \(C_0\) makes the error in expectation
negligible. Multiplication by \(A\in[0,1]\) does not alter this bound.
Thus this first precision choice precedes all moment constructions.

Orthogonalize the Riesz vectors of the retained physical sine tests.
The subspace and its real orthogonalization are independent of
\(M,\nu\). Let \(Y\in\R^d\) be the vector of resulting tests on
\(Q\). The truncated exponential is
\[
 g(Y)=\exp(-c_1-l^{\mathsf t}Y-\tfrac12Y^{\mathsf t}DY),
 \qquad I+\operatorname{Re}D\succeq\delta I,
 \qquad \delta=L^{-r_0},
\]
with \(D\) complex symmetric, \(d,\|D\|,\|l\|\) polynomial, and
\(\operatorname{Re}c_1\ge-L^C\). Operator-norm convergence of the
truncated kernel preserves this gap. Each coordinate of \(Y\) is a
centered sum of a physical cosine polynomial \(f(s)\), with
polynomial degree and coefficients. The inverse-polynomial Gram
bound established in Lemma~\ref{lem:04-quadratic} justifies this last
assertion after orthogonalization.

\begin{lemma}[Weighted polynomial approximation]\label{lem:04-hermite}
For the preceding Gaussian exponential, choose \(a=1-\delta/2\).
Given any fixed precision exponent, there is a polynomial \(P(Y)\)
of polynomial total degree \(K_1\) such that
\[
 \sup_{y\in\R^d}e^{-a|y|^2/2}|g(y)-P(y)|\le e^{-L^{B_1}}.
\]
Its coefficients are holomorphic in \(M\) and bounded by an
exponential of a fixed polynomial. The degree and coefficient bounds
are fixed before any endpoint expansion.
\end{lemma}

\begin{proof}
Use the Gaussian probability measure
\(d\mu_a=(a/\pi)^{d/2}e^{-a|y|^2}dy\), and its orthonormal Hermite
polynomials defined by
\[
 e^{\sqrt{2a}z^{\mathsf t}y-z^{\mathsf t}z/2}
       =\sum_\alpha H_\alpha(y)\frac{z^\alpha}{\sqrt{\alpha!}}.
\]
Gaussian integration proves orthogonality. For completeness,
polynomials are dense: an orthogonal complement would have an entire
Gaussian-weighted Laplace transform with every derivative at zero
equal to zero; Fourier uniqueness then makes the complement zero.
The integrability needed here follows from Gaussian decay.

Let \(c_\alpha=\int gH_\alpha\,d\mu_a\). Completing the square gives
\begin{align*}
 \sum_\alpha c_\alpha\frac{z^\alpha}{\sqrt{\alpha!}}
 &=e^{-c_1}\det(I+D/(2a))^{-1/2}\\
 &\quad\times\exp\!\left\{
 \tfrac12l^{\mathsf t}(2aI+D)^{-1}l
 -\sqrt{2a}z^{\mathsf t}(2aI+D)^{-1}l
 -\tfrac12z^{\mathsf t}D(2aI+D)^{-1}z\right\}.
\end{align*}
The square root is the analytic branch specified by the Gaussian
integral. Accretivity ensures its existence and polynomial inverse
norms. Set \(B=D(2aI+D)^{-1}\). The identity
\[
 (2aI+D)^*(2aI+D)-D^*D
       =4a(aI+\operatorname{Re}D)\succeq2a\delta I
\]
shows \(\|B\|\le1-L^{-C}\). The other factors in the generating
function have size at most \(\exp(L^C)\). Therefore, for some
\(r=1+L^{-C_1}>1\), integrating the squared generating function at
\(rz\) against \(\pi^{-d}e^{-|z|^2}\,dz\) on \(\C^d\) gives
\[
 \sum_\alpha r^{2|\alpha|}|c_\alpha|^2\le e^{L^C}.
\]
This follows from monomial orthogonality; convergence follows from
\(r^2\|B\|<1\), with inverse-polynomial margin.

The Hermite generating identity also gives Mehler's diagonal formula
\cite[equation~18.18.28]{DLMF18}
\[
 \sum_\alpha t^{|\alpha|}|H_\alpha(y)|^2
 =(1-t^2)^{-d/2}
       \exp\!\left(\frac{2at}{1+t}|y|^2\right),\quad 0<t<1.
\]
One direct derivation writes each one-variable unnormalized Hermite
polynomial as the moment of \(\sqrt{2a}y+iG\), for a standard real
Gaussian \(G\), and integrates the summed exponential for two
independent copies. Taking \(t=r^{-1}\) bounds the right side by
\((1-r^{-2})^{-d/2}e^{a|y|^2}\). Cauchy--Schwarz now bounds the
Hermite tail above degree \(K\), after multiplying by
\(e^{-a|y|^2/2}\), by
\(e^{L^C}r^{-K/2}\). The prefactor includes the polynomial-dimensional
Mehler factor. A polynomially large \(K=K_1\) attains any specified
\(B_1\). Completeness and local uniform convergence identify the
series with \(g\). Expanding the retained Hermite polynomials in
monomials gives the asserted coefficient bound, since the dimension
and degree are polynomial. The Gaussian coefficient integrals,
uniformly dominated on compact \(M\)-sets, prove holomorphy.
\end{proof}

Gaussian linearization and \eqref{eq:G-13} imply
\[
 \E_0e^{a|Y|^2/2}\le e^{CL}(1-a)^{-d/2}.
\]
Consequently the weighted error in Lemma~\ref{lem:04-hermite} gives
an error in expectation smaller than \(e^{-L^b}\) by fixing \(B_1\)
after \(d,r_0\). From this point onward \(K_1\) and all polynomial
coefficient bounds are fixed.

We next replace each constituent \(f(s(t))\) by
\[
 f(s(t))=g_f(\cos t)+e_f(t)+r_f(t),\qquad
 \deg g_f\le d_g=\lfloor m/L^{D_1}\rfloor,
 \qquad \|r_f\|_\infty\le e^{-L^{B_2}},
\]
where \(e_f\) is polynomially bounded and supported on
\(\min(t,\pi-t)\le L^{D_E}/m\). Choose \(D_1\) first so that
\(K_1d_g=o(m)\), with ample fixed margin, then choose \(B_2\) after
the degree and coefficient bounds of \(P\), and finally choose
\(D_E\).
To obtain this decomposition, multiply \(f(s(t))\) by an even-reflected
Gevrey mask which is zero up to distance \(L^{D_E}/(2m)\) from an
endpoint and one beyond \(L^{D_E}/m\). The omitted part is \(e_f\).
Outside the omitted region, integration of
\(\lambda_L(s)\asymp e^{-\gamma|s|}\) places \(s(t)\) at a growing
distance from the folded endpoint buffer. There \(\theta'\asymp
\min(t,\pi-t)\), with relatively bounded analytic derivatives on
fixed small \(s\)-disks. The analytic inverse theorem therefore
extends \(s(t)\) to a disk of radius
\(c\min(t,\pi-t)\), with bounded displacement. Shrinking this radius
by a polynomial to handle the physical sine bandwidth gives masked
derivative bounds
\[
 L^C\big(C'm/L^{D_E-C_2}\big)^k(k!)^{P'},
\]
with fixed \(C_2,P'\) independent of sufficiently large \(D_E\).
The resulting cosine Fourier tail above \(d_g\) is at most a
polynomial times
\(e^{O(L)}\exp[-cL^{(D_E-C_2-D_1)/P'}]\). Increasing \(D_E\)
proves the required sup error. Both the real sup norm of \(g_f\)
and its individual cosine coefficients are polynomially bounded.

Keep the original centering constants of \(Y\). The uncentered sums
of all retained tests and edge functions have magnitude at most
\(e^{CL}L^C\). A polynomial of total degree \(K_1\), with the
coefficient bounds already fixed, has derivative at most
\(e^{L^{C'}}\) on this range. The total perturbation of its arguments
is at most \(e^{CL-L^{B_2}}\). Thus the mean-value formula bounds
the error in \(P\) uniformly by
\(e^{L^{C'}+CL-L^{B_2}}\), which was made negligible when choosing
\(B_2\). This estimate is completed before any subset of endpoint
particles is introduced.

\subsection{A finite endpoint-subset expansion}

Set \(x=\cos t\). For each label the normalized reference law is the
Legendre projection process for Lebesgue measure on \([-1,1]\).
Let \(E_L\) be the two endpoint intervals corresponding to
\(\min(t,\pi-t)\le L^{D_E}/m\); their total length is
\(O(L^{2D_E}/m^2)\). Enlarge \(D_E\), if necessary, so that every
factor of \(A\) differs from one only on \(E_L\). This is possible
because those factors require physical distance bounded from a wall.
They are bounded in \([0,1]\) and independent of \(M,\nu\).

Let \(S_1\) denote the vector of global sums of the polynomials
\(g_f(x)\), and let \(B\) be the actual labeled endpoint-particle
set. At fixed \(S_1\) and centerings, define \(H(D)\), for
\(D\subseteq B\), by using only points of \(D\) in the edge sums
and in the product of endpoint factors of \(A\). Define
\[
 I(B')=\sum_{D\subseteq B'}(-1)^{|B'|-|D|}H(D).
\]
M\"obius inversion on subsets gives the exact identity
\(H(B)=\sum_{B'\subseteq B}I(B')\): the coefficient of \(H(D)\)
in the sum is \((1-1)^{|B\setminus D|}\). Each \(I(B')\) is a
polynomial in \(S_1\) of total degree at most \(K_1\). Uniformly on
integer configurations,
\[
 |I(B')|\le 2^{|B'|}e^{L^{C_2'}}.
\]
The exponent here is independent of \(|B'|\), since every possible
subset contains at most \(O(m)\) points and the polynomial degree
has already been fixed.

For the normalized Legendre polynomials
\(p_i=\sqrt{(2i+1)/2}\,P_i\), one has \(|P_i(x)|\le1\) on the
real interval. A direct verification uses
\[
 P_i(x)=\frac1\pi\int_0^\pi
       \big(x+i\sqrt{1-x^2}\cos\phi\big)^i\,d\phi.
\]
Summing the geometric series proves its generating function
\((1-2xz+z^2)^{-1/2}\), hence its equality with the polynomials from
the Legendre recurrence. The integrand has modulus at most one.
Therefore \(\Pi_n(x,x)=\sum_{i<n}p_i(x)^2\le n^2/2\).
The factorial densities are minors of this positive kernel; the
Gram determinant bound gives, for the total endpoint count \(N_E\)
of both labels,
\[
 \E_0\binom{N_E}{k}\le\frac{\beta_L^k}{k!},\qquad
 \beta_L=C L^{2D_E}.
\]
For two labels this follows either from the direct-sum kernel or by
summing the binomial convolution of their separate bounds.
Thus truncating the subset expansion at \(|B'|\le K_e\) has error
at most
\[
 e^{L^{C_2'}}\sum_{k>K_e}\frac{(2\beta_L)^k}{k!}.
\]
If \(K_e+1\ge4e\beta_L\), this tail is at most
\(2e^{L^{C_2'}}(2e\beta_L/(K_e+1))^{K_e+1}\).
Choose the polynomial \(K_e\) large enough to attain the desired
precision. Crucially, this does not increase \(K_1\): the endpoint
subset size and the global polynomial degree are separate orders.

\subsection{Continuation of the retained factorial moments}

The determinant formulas here are the standard consequences of
Andr\'eief's identity for an orthogonal-polynomial ensemble
\cite[equations (1.3), (3.1), and Section 5]{TracyWidom1998}.
We give the expansion because the subsequent
continuation uses its precise finite degrees.

We describe one label, with integer dimension \(n\), later continued
to \(\nu\). Introduce formal source variables \(\mathbf t=(t_j)\),
and \(G(x)=\sum_jt_jg_j(x)\), where the \(g_j\) are the global
polynomial tests. The order-\(k\) factorial density, tilted by
\(e^{\sum_iG(x_i)}\) and without the divisor \(k!\), is
\[
 \mathcal D_n(\mathbf t)
       \det[\mathcal C_n(x_h,x_{h'};\mathbf t)]_{h,h'\le k},
\]
where
\[
 \mathcal D_n=\det_{\Pi_n}(\Pi_ne^G\Pi_n),\qquad
 \mathcal C_n=\Pi_n(\Pi_ne^G\Pi_n|_{\Pi_n})^{-1}\Pi_ne^G.
\]
These identities are formal at \(\mathbf t=0\), where the inverse
exists. To prove them, insert an additional factor
\(\prod_i(1+h(x_i))\) in the tilted expectation. Integration of the
two polynomial evaluation determinants gives the Gram determinant
with weight \(e^G(1+h)\). Factor out its value at \(h=0\), and
expand the remaining determinant in its minors. The coefficient of
\(h(x_1)\cdots h(x_k)\) is exactly the formula above.

For a labeled endpoint set write
\[
 I(B')=\sum_{|\alpha|\le K_1}
        c_\alpha(B';M,\nu)S_1^\alpha.
\]
The coefficients are holomorphic in \(M,\nu\): they are obtained
from \(P\), centerings linear in \(\nu\), and sums or products of
bounded real endpoint functions evaluated at the specified points.
Their magnitudes are at most \(e^{L^C}\). For each split of
\(|B'|=k\) into \(k_+,k_-\), integrate these coefficients against
\(\partial_{\mathbf t}^{\alpha}|_{\mathbf t=0}\) of the product
of the two tilted factorial densities, and divide by \(k_+!k_-!\).
This is precisely the expected subset sum. Only source degrees at
most \(K_1\) are required, regardless of \(k\le K_e\).

In the polynomial-index basis, multiplication by \(x\) is the
tridiagonal Jacobi matrix \(J\), with off-diagonal coefficients
\[
 a_i=\frac{i}{\sqrt{(2i-1)(2i+1)}}
       =(4-i^{-2})^{-1/2}.
\]
At source degree \(r\), multiplication by \(e^G\) has bandwidth at
most \(rd_g\). In
\[
 \log\mathcal D_n
   =\operatorname{Tr}_{\Pi_n}
         \log(I+\Pi_n(e^G-1)\Pi_n),
\]
a closed index path of degree \(r\) has total displacement at most
\(rd_g\). If its starting index lies farther than \(rd_g\) below
\(n\), it never meets the upper projection boundary; every inserted
\(\Pi_n\) can be removed. The complete coefficient then equals
that of \(\log e^G=G\), and vanishes for \(r\ge2\). This argument
is an identity of formal series, including the cancellations among
the finitely many logarithm terms. All higher cumulants therefore
involve only indices within \(2K_1d_g+O(1)\) of \(n\).

The kernel correction has the same property. Exactly,
\(\mathcal C_n\Pi_n=\Pi_n\) and
\((I-\Pi_n)\mathcal C_n=0\), so
\(\mathcal C_n-\Pi_n\) has rows below \(n\) and columns at least
\(n\). Its degree-\(r\) inverse-series coefficient has bandwidth
\(rd_g\), forcing both indices within \(rd_g\) of \(n\).
All these assertions concern polynomial kernels coefficient by
coefficient; pointwise evaluation at endpoints is legitimate.

Continue the required finite blocks by replacing \(a_{n+j}\) with
\[
 a_{\nu+j}=(4-(\nu+j)^{-2})^{-1/2},
\]
using the branch near \(1/2\), on \(|\nu-m|\le cm\) for a small
fixed \(c\). Choose a fixed offset block, for example
\(|j|\le m/8\), and let the projection test the sign of the offset.
Our choice \(K_1d_g=o(m)\) ensures that no retained coefficient
reaches its artificial endpoints. The block differs in norm by
\(O(m^{-2})\) from the real constant tridiagonal matrix with entries
\(1/2\). The latter has spectrum in \([-1,1]\). On a contour at
a sufficiently large constant times \(m^{-2}\) from this interval,
a resolvent identity bounds the continued resolvent by \(O(m^2)\).
For the Chebyshev polynomial \(T_p(x)=\cos(p\arccos x)\), the
quadratic-root formula gives
\( |T_p(x)|\le C e^{Cp/m}\) on this contour. Expanding each
\(g_j\) in its cosine coefficients and applying the resolvent
formula thus bounds \(\|g_j(J)\|\) by \(e^{CL}L^C\).
There are at most exponentially many indices but only polynomially
many factors in any retained term; all products, traces, and formal
coefficients therefore have bounds \(e^{L^C}\). This resolvent
estimate is essential because an absolute enumeration of paths of
length \(d_g\) would give a much larger bound.

The degree-one trace
\(T_j(n)=\operatorname{Tr}(\Pi_ng_j\Pi_n)\) requires a separate
construction. Its forward difference
\(T_j(n+1)-T_j(n)=(g_j(J))_{nn}\) is a near-index entry, hence has
an analytic continuation \(h_j(\nu)\) on \(|\nu-m|\le cm\)
with bound \(e^{CL}L^C\). If
\(h_j(m+w)=\sum_{p\ge0}h_{jp}w^p\), Cauchy's estimate gives
\(|h_{jp}|\le e^{CL}L^C(C/m)^p\). Let \(S_p\) be the polynomial
with \(S_p(0)=0\) and
\(S_p(w+1)-S_p(w)=w^p\). Define
\[
 \widehat T_j(m+w)=T_j(m)+\sum_{p=0}^{P_3}h_{jp}S_p(w).
\]
This agrees, at integer \(w\), with summing the truncated forward
differences in either direction from zero. The generating function
\[
 \sum_{p\ge0}S_p(w)\frac{t^p}{p!}
          =\frac{e^{tw}-1}{e^t-1}
\]
and Cauchy's estimate on a circle of radius comparable to
\((p+1)/(p+1+|w|)\) give
\( |S_p(w)|\le[C(1+|w|+p)]^{p+1}\).
For \(|w|\le\sqrt m\) and any fixed polynomial \(P_3\), the
continued trace is consequently bounded by \(e^{CL}L^C\), uniformly
in the chosen truncation power for sufficiently large \(L\).
Its integer-grid error is at most
\(e^{CL}L^C\sqrt m(C/\sqrt m)^{P_3+1}\).
Increasing the polynomial order gives any required
\(e^{-L^{B_3}}\) precision.

The same construction continues the reference masses. The exact norm
product gives
\[
 \Delta^2\log J_n(w_0)
   =\log\!\left(1+\frac1{(2n+1)(2n+3)}\right).
\]
Keep the exact \(\log J_m\) and \(\Delta\log J_m=O(m^{-1})\),
Taylor-expand this analytic second difference at \(m\), and sum
twice by polynomial antidifferences. Its Taylor coefficients are
\(O(m^{-2})(C/m)^p\). The polynomial twice-antidifference with
values zero at \(w=0,1\) has generating function
\[
 \frac{e^{tw}-1-w(e^t-1)}{(e^t-1)^2}
\]
and bound \([C(1+|w|+p)]^{p+2}\). The resulting analytic logarithm
stays \(O(L)\) on \(|w|\le\sqrt m\). Integer errors acquire at
most a factor \(m\) when summed twice; after exponentiation they
still attain arbitrary \(e^{-L^{B_3}}\) accuracy. These uniform
size bounds do not grow with the final polynomial Taylor order.

Endpoint evaluations are continued directly. The constant kernel is
represented by Christoffel--Darboux~\cite[equation~18.2.12]{DLMF18}:
\[
 \Pi_n(x,y)=a_n
   \frac{p_n(x)p_{n-1}(y)-p_{n-1}(x)p_n(y)}{x-y},
\]
with the removable diagonal filled by its limit. For \(x\) near
\(+1\), continue a wavefunction with index \(z=\nu+j\) by
\[
 p_z(x)=\sqrt{(2z+1)/2}
 \sum_{r=0}^{\infty}
 \frac{\prod_{l=0}^{r-1}(z-l)(z+l+1)}{(r!)^2}
                  \left(\frac{x-1}{2}\right)^r.
\]
This is the Rodrigues expansion at integer \(z\). For
\(|j|\le m/8\), \(|\nu-m|\le\sqrt m\), and
\(|x-1|\le C L^{2D_E}/m^2\), the terms with \(r\le4m\) are
bounded by \((C L^{D_E})^{2r}/(r!)^2\), apart from the square-root
prefactor. For \(r>4m\), the ratio of successive absolute
majorants is at most \(C|x-1|<1/2\). Thus the function is bounded
by \(e^{L^C}\), also on slightly larger complex endpoint disks.
Cauchy estimates give the same type of bound for its first derivative,
including the factor \(m^2\) from the disk radius. This controls
same-endpoint diagonal quotients; opposite-endpoint denominators
are bounded away from zero.

Near \(-1\), use the value at \(-x\) and multiply by
\((-1)^j\), omitting any putative \((-1)^\nu\). At integer \(n\)
this prescription gives
\(\widetilde p_{n+j}(x)=d_n(x)p_{n+j}(x)\), where
\(d_n(x)=1\) near \(+1\) and \(d_n(x)=(-1)^n\) near \(-1\).
Consequently both the Christoffel--Darboux kernel (with its original
\(x-y\) denominator) and every retained correction kernel are
changed by \(K(x,y)\mapsto d_n(x)K(x,y)d_n(y)\). For any endpoint
set their determinant is unchanged, because the row and column
multipliers contribute \(\prod_h d_n(x_h)^2=1\). This proves
agreement at integer dimensions even for configurations using both
endpoints, while providing a holomorphic continuation in \(\nu\).

\begin{proof}[Proof of Proposition~\ref{prop:04-interpolation}]
Use the polynomial and endpoint-subset approximations just constructed.
In each retained subset integral, replace the higher cumulants and
kernel corrections by their continued near-index formulas, the
constant kernel by the continued Christoffel--Darboux expression,
and the linear traces and masses by their polynomial-antidifference
approximants. Reconstruct \(\mathcal D_n\) by formal exponentiation
through source degree \(K_1\), and take the indicated source
derivatives. The two labels are combined by independence.

Every expression before endpoint integration is holomorphic in
\(\nu,M\), with a uniform bound \(e^{L^C}\). Indeed source types,
source orders, and determinant sizes are polynomial. Near-index
blocks have length \(O(m)\), so an expression with polynomially
many index sums has at most \(\exp(\operatorname{poly}(L))\)
terms. Permutation expansions, source coefficient extractions, and
factorials of polynomial orders satisfy the same bound. Coefficients
\(c_\alpha\) include only the fixed polynomial \(P\), centerings,
and bounded endpoint data. Integration over fixed real endpoint
regions preserves these bounds and holomorphy by dominated
convergence (or Morera's theorem). No analyticity of the endpoint
factors themselves as functions of their real positions is needed.

We record the order of the error choices to exclude a circular
precision requirement. First fix \(B_0\) using the original
integer-law expectation bound. This fixes the dimension and gap in
the Gaussian approximation. Next fix \(B_1\) and \(K_1\); then
\(D_1,B_2,D_E\), in that order; finally choose \(K_e\) for the
factorial tail. All errors to this point have already been bounded
on the actual integer laws by \(e^{-L^b}\), with a freely reserved
margin. They are not passed through an analytic continuation formula.
The remaining formulas, through degree \(K_1\), are polynomial in
the linear trace coefficients and in the kernel entries. Their
sensitivity to the final trace and normalization errors is at most
\(e^{L^{C_*}}\), for an exponent \(C_*\) depending only on the
choices already made. The preceding size bounds for trace and mass
approximants are independent of their final polynomial Taylor order.
Choose \(B_3>\max(b,C_*)+1\), and the corresponding orders, last.
The resulting errors are negligible after this final amplification.
Taking the initially reserved margins, for example replacing each
requested error by \(e^{-2L^b}\), makes their finite sum at most
\(e^{-L^b}\) for large \(L\). This proves both assertions in
\eqref{eq:G-20}.
\end{proof}

\begin{lemma}[Transfer from the integer grid]\label{lem:04-grid}
Suppose \(M\) is in the fixed polynomial box and
\(|F_\nu(M)|\le B\) at all integer grid points of
\(|\nu_\sigma-m|\le\sqrt m\), with
\(B\ge e^{-L^{b'}}\), \(b'<b\). If \(\nu^*-m\) has polynomial
size, then
\begin{equation}
 |\widehat F(\nu^*,M)|\le B\exp(o(L)).
 \label{eq:G-21}
\end{equation}
The estimate is uniform under the stated bounds; in particular it
applies to \(\nu^*=m'(M)\).
\end{lemma}

\begin{proof}
Taylor-expand \(\widehat F\) in its two dimension variables at
\((m,m)\). Its bound \(e^{L^C}\) on the radius-\(\sqrt m\)
polydisk implies that a polynomial \(P\) of polynomial degree
\(K\) approximates it on the half-radius polydisk to an error
\(o(B)\), by choosing \(K\) so that \(e^{L^C}2^{-K}=o(B)\).
Put \(R=\lfloor\sqrt m/3\rfloor\), and let \(S\) be the sup norm
of \(P\) on the real square \([m-R,m+R]^2\).
The rescaled one-variable Chebyshev expansion has coefficients at
most twice its sup norm and \(\|T_k'\|_{[-1,1]}\le k^2\).
Holding the other variable real therefore bounds each partial
derivative by \(CK^3S/R\). Comparing a maximizer with a nearest
integer point gives
\[
 S\le \max_{\text{integer square}}|P|+CK^3S/R.
\]
Since \(R\) is exponential and \(K\) polynomial, the second
coefficient is \(o(1)\). Proposition~\ref{prop:04-interpolation}
and \(b>b'\) then give \(S\le(1+o(1))B\).

The tensor Chebyshev coefficients are at most \(4S\), and there
are at most \((K+1)^2\) of them. At the target, each rescaled
coordinate has magnitude \(\operatorname{poly}(L)/R\). Its
complex inverse-cosine angle differs from \(\pi/2\) by the same
order, so every Chebyshev basis value up to degree \(K\) is
bounded by \(\exp(CK\operatorname{poly}(L)/R)=1+o(1)\).
Thus \(|P(\nu^*)|\le C(K+1)^2B\), which is
\(B\exp(o(L))\). The Taylor error is negligible relative to \(B\).
\end{proof}

For fixed explicit data, the prescribed totals \(m'_\sigma(M)\)
are affine in \(M\). Under the retained polynomial restrictions,
\(m'(M)-m\) has polynomial size and lies in the dimension polydisk.
Consequently
\[
 H(M)=\widehat F(m'(M),M)
\]
is holomorphic throughout the chosen polynomial \(M\)-disk. At every
valid integer \(M\) in that disk, it approximates the original
integer-dimensional integral to the precision in \eqref{eq:G-20}.
The grid-transfer lemma supplies the sharper bounds on \(H\) away
from those integer values; they do not follow from the coarse
holomorphic bound alone.

